\chapter{Integrals with Arithmetic Inside}

\begin{center}\itshape
The derivation of these properties requires no more than some second year calculus.\\
\upshape --- Frits Beukers, on his proof that $\zeta(3)$ is irrational, in his farewell lecture \emph{A passion for numbers} (Utrecht, 2019)
\end{center}

In June 1978 Roger Apéry, then sixty-one, stood up at the Journées Arithmétiques in Marseille and announced that
$\zeta(3)=1+\frac1{2^3}+\frac1{3^3}+\cdots$ is irrational. Euler had computed $\zeta(2)=\frac{\pi^2}6$ in 1735, and nothing at all was known
about the odd values; a proof seemed out of reach. The lecture did not help. Alf van der Poorten, who was in the room, later
wrote in the \emph{Mathematical Intelligencer}: ``Though there had been earlier rumours of his claiming a proof, scepticism was
general. The lecture tended to strengthen this view to rank disbelief.'' Apéry wrote formula after formula on the board, such as
\[
\zeta(3)=\frac52\sum_{n\ge1}\frac{(-1)^{n-1}}{n^3\binom{2n}n},
\]
with little explanation of where they came from. But some listeners, Henri Cohen, Hendrik Lenstra and van der Poorten among them,
noticed that the formulas could be checked, and they checked. Cohen, with help from Don Zagier, filled in the missing steps, and
the complete proof was presented at the International Congress of Mathematicians in Helsinki that August.

In the audience in Marseille sat a young Dutch doctoral student at his first international conference, Frits Beukers. A few months
later, as he tells it, he found a proof ``inspired by Apéry's, but otherwise completely different and much simpler''. It rests on a
single triple integral,
\[
I_n=\int_0^1\!\!\int_0^1\!\!\int_0^1\frac{\bigl(x(1-x)y(1-y)z(1-z)\bigr)^n}{\bigl(1-(1-xy)z\bigr)^{n+1}}\dd x\dd y\dd z,
\]
which has two properties: $I_n=b_n\zeta(3)-a_n$ for integers $a_n$, $b_n$ divided by a known denominator, and $I_n$ is positive but
astonishingly small. The proof appeared in 1979 and, in Beukers's words, went ``viral''. You will rebuild it in
Problem~\ref{pb:beukersthree}.

Twenty years later, Tanguy Rivoal (2000) proved that infinitely many of the numbers $\zeta(3),\zeta(5),\zeta(7),\dots$ are irrational,
and Wadim Zudilin (2001) that at least one of $\zeta(5),\zeta(7),\zeta(9),\zeta(11)$ is. Nobody knows which one. Whether $\zeta(5)$ is
irrational, or Euler's constant $\gamma$, is still open. This chapter is about integrals that carry arithmetic inside them: small
positive numbers that are secretly integer combinations of constants, sums of binomial coefficients that are areas, fractional
parts whose average is $1-\gamma$, and sums of uniform random numbers that count permutations.

\begin{think}
How can an integral prove that a number $\alpha$ is irrational? Suppose you have integers $a_n,b_n$ with $b_n\alpha-a_n\ne0$ for every
$n$ and $b_n\alpha-a_n\to0$. Why is that enough? Then try to see why a positive integral that tends to $0$ is a natural candidate for
$b_n\alpha-a_n$.
\end{think}

\section{Small integrals, irrational numbers}

The answer to the think box is the engine of the whole section: if $\alpha=\frac pq$, then $q(b_n\alpha-a_n)=pb_n-qa_n$ is a non-zero
integer, so its absolute value is at least $1$, and it cannot tend to $0$. Integrals are a good source of such numbers because an
integrand can be made small everywhere while its expansion produces only rational numbers and the constant $\alpha$. Throughout,
$d_n=\operatorname{lcm}(1,2,\dots,n)$, and
\[
P_n(x)=\frac1{n!}\,\frac{\mathrm d^n}{\mathrm dx^n}\bigl(x^n(1-x)^n\bigr)
\]
is the $n$-th shifted Legendre polynomial.

\begin{exercise}[$\star\star$]\label{pb:logtwo}
\emph{The irrationality of $\log2$.}
\begin{enumerate}[label=(\alph*)]
\item Show that $P_n$ has integer coefficients, and that
\[
I_n=\int_0^1\frac{P_n(x)}{1+x}\dd x=\int_0^1\frac{x^n(1-x)^n}{(1+x)^{n+1}}\dd x.
\]
\item Show that $d_nI_n=a_n+b_n\log2$ with $a_n,b_n\in\Z$.
\item Show that $0<I_n\le(3-2\sqrt2)^n$, and deduce that $\log2$ is irrational.
\end{enumerate}
\end{exercise}

\begin{exercise}[$\star\star$]\label{pb:hermite}
\emph{Powers of $e$.}
\begin{enumerate}[label=(\alph*)]
\item Show that $e^k$ is irrational for every integer $k\ge1$.
\item Deduce that $e^r$ is irrational for every rational $r\ne0$, and that $\log r$ is irrational for every rational $r>0$, $r\ne1$.
\end{enumerate}
\end{exercise}

\begin{exercise}[$\star$]\label{pb:doubleint}
\emph{Double integrals that know $\zeta(2)$ and $\zeta(3)$.} Let $r,s\ge0$ be integers.
\begin{enumerate}[label=(\alph*)]
\item Show that $\displaystyle\int_0^1\!\!\int_0^1\frac{x^ry^s}{1-xy}\dd x\dd y$ equals $\frac1{r-s}\bigl(\frac1{s+1}+\cdots+\frac1r\bigr)$ if $r>s$, and
$\zeta(2)-\sum_{k=1}^r\frac1{k^2}$ if $r=s$.
\item Show that $\displaystyle\int_0^1\!\!\int_0^1\frac{-\log(xy)}{1-xy}\,x^ry^s\dd x\dd y$ equals $\frac1{r-s}\bigl(\frac1{(s+1)^2}+\cdots+\frac1{r^2}\bigr)$ if $r>s$, and
$2\bigl(\zeta(3)-\sum_{k=1}^r\frac1{k^3}\bigr)$ if $r=s$.
\item Show that $\displaystyle\int_0^1\!\!\int_0^1\frac{\mathrm dx\dd y}{1+xy}=\frac{\pi^2}{12}$.
\end{enumerate}
\end{exercise}

\begin{exercise}[$\star\star\star$]\label{pb:beukerstwo}
\emph{Beukers's proof for $\zeta(2)$.} Let
\[
J_n=\int_0^1\!\!\int_0^1\frac{(1-y)^nP_n(x)}{1-xy}\dd x\dd y.
\]
\begin{enumerate}[label=(\alph*)]
\item Show that $d_n^2J_n=a_n+b_n\zeta(2)$ with $a_n,b_n\in\Z$.
\item Show that
\[
J_n=(-1)^n\int_0^1\!\!\int_0^1\frac{x^n(1-x)^ny^n(1-y)^n}{(1-xy)^{n+1}}\dd x\dd y.
\]
\item Show that $\frac{x(1-x)y(1-y)}{1-xy}\le\bigl(\frac{\sqrt5-1}2\bigr)^5$ on $[0,1]^2\setminus\{(1,1)\}$, and deduce that $0<|J_n|\le\zeta(2)\bigl(\frac{\sqrt5-1}2\bigr)^{5n}$.
\item Deduce that $\zeta(2)$ is irrational. (This gives a second proof that $\pi^2$ is irrational.)
\end{enumerate}
\end{exercise}

\begin{exercise}[Journey]\label{pb:beukersthree}
\emph{Beukers's proof that $\zeta(3)$ is irrational.} Let
\[
K_n=\int_0^1\!\!\int_0^1\frac{-\log(xy)}{1-xy}\,P_n(x)P_n(y)\dd x\dd y.
\]
\begin{enumerate}[label=(\alph*)]
\item Show that $d_n^3K_n=a_n+b_n\zeta(3)$ with $a_n,b_n\in\Z$.
\item Show that $\frac{-\log(xy)}{1-xy}=\int_0^1\frac{\mathrm dz}{1-(1-xy)z}$ for $0<xy<1$.
\item Integrating by parts $n$ times in $x$, then substituting $w=\frac{1-z}{1-(1-xy)z}$, then integrating by parts $n$ times in $y$, show that
\[
K_n=\int_0^1\!\!\int_0^1\!\!\int_0^1\frac{x^n(1-x)^ny^n(1-y)^nw^n(1-w)^n}{\bigl(1-(1-xy)w\bigr)^{n+1}}\dd x\dd y\dd w.
\]
\item Show that $\frac{x(1-x)y(1-y)w(1-w)}{1-(1-xy)w}\le(\sqrt2-1)^4$ on $[0,1]^3$ (where defined), and that $0<K_n\le2\zeta(3)(\sqrt2-1)^{4n}$.
\item Using the prime number theorem in the form $\log d_n\sim n$, conclude that $\zeta(3)$ is irrational.
\end{enumerate}
\end{exercise}

\begin{figure}[ht]
\centering
\fig{ch09-beukers}
\caption{The engine of Beukers's proofs at work: $\log_{10}$ of $d_n^2|J_n|$ (Problem~\ref{pb:beukerstwo}) and of $d_n^3|K_n|$
(Problem~\ref{pb:beukersthree}) for $n\le40$, computed exactly from the integer coefficients. Both are non-zero integer combinations of
$1$ and $\zeta(2)$, or of $1$ and $\zeta(3)$, and both tend to $0$ geometrically. Since $d_n=e^{n(1+o(1))}$, the eventual rates are $e^2\bigl(\frac{\sqrt5-1}2\bigr)^5=0.67$ and
$e^3(\sqrt2-1)^4=0.59$ per step; the proofs only need the cruder bound $d_n\le3^n$, which gives $0.81$ and $0.79$.}
\label{fig:beukers}
\end{figure}

\begin{exercise}[$\star\star$]\label{pb:sondow}
\emph{Euler's constant as a double integral (Sondow, 2005).}
\begin{enumerate}[label=(\alph*)]
\item Show that
\[
\gamma=\int_0^1\!\!\int_0^1\frac{1-x}{(1-xy)\bigl(-\log(xy)\bigr)}\dd x\dd y.
\]
\item Show similarly that replacing $1-xy$ by $1+xy$ gives $\log\frac4\pi$.
\end{enumerate}
\end{exercise}

\begin{exercise}[$\star$]\label{pb:dalzell}
\emph{Why $\pi<\frac{22}7$.}
\begin{enumerate}[label=(\alph*)]
\item Show that $\displaystyle\int_0^1\frac{x^4(1-x)^4}{1+x^2}\dd x=\frac{22}7-\pi$, and deduce that $\frac1{1260}<\frac{22}7-\pi<\frac1{630}$.
\item Show that for every $n\ge1$ there is a rational $r_n$ with $\displaystyle\int_0^1\frac{x^{4n}(1-x)^{4n}}{1+x^2}\dd x=r_n+(-4)^n\frac\pi4$, and that
$\bigl|\pi+4(-4)^{-n}r_n\bigr|<4^{1-5n}$.
\end{enumerate}
\end{exercise}

\section{Central binomial coefficients}

The numbers $\binom{2n}n$ are everywhere in Part~I: they count lattice paths and returns of random walks (Tome~3), they bound
primes (Tome~1), and they are Beta integrals in disguise (Tome~4). Turning $\frac1{\binom{2n}n}$ into an integral, or recognising a
series of them as the Taylor series of an inverse trigonometric function, sums series that look hopeless at first sight.

\begin{exercise}[$\star\star$]\label{pb:centralsum}
\emph{Reciprocals of central binomial coefficients.}
\begin{enumerate}[label=(\alph*)]
\item Show that $\displaystyle\frac1{n\binom{2n}n}=\int_0^1t^{n-1}(1-t)^n\dd t$ for $n\ge1$, and deduce that $\displaystyle\sum_{n\ge1}\frac1{n\binom{2n}n}=\frac\pi{3\sqrt3}$.
\item Show that $\displaystyle\sum_{n\ge0}\frac1{\binom{2n}n}=\frac43+\frac{2\pi\sqrt3}{27}$.
\end{enumerate}
\end{exercise}

\begin{exercise}[$\star\star$]\label{pb:arcsin}
\emph{The square of the arcsine.}
\begin{enumerate}[label=(\alph*)]
\item Show that $y=\arcsin^2x$ satisfies $(1-x^2)y''-xy'=2$ on $]-1,1[$, with $y(0)=y'(0)=0$.
\item Deduce that for $|x|\le1$
\[
\arcsin^2x=\frac12\sum_{n\ge1}\frac{(2x)^{2n}}{n^2\binom{2n}n}.
\]
\item Deduce that $\displaystyle\sum_{n\ge1}\frac1{n^2\binom{2n}n}=\frac{\pi^2}{18}$ and $\displaystyle\sum_{n\ge1}\frac{4^n}{n^2\binom{2n}n}=\frac{\pi^2}2$.
\end{enumerate}
\end{exercise}

\begin{exercise}[$\star$]\label{pb:lehmer}
\emph{Series with $\binom{2n}n/4^n$.} Let $c_n=\binom{2n}n4^{-n}$.
\begin{enumerate}[label=(\alph*)]
\item Show that $\displaystyle\sum_{n\ge0}\frac{c_n}{2n+1}=\frac\pi2$ and $\displaystyle\sum_{n\ge1}\frac{c_n}n=2\log2$.
\item Show that $\sum_nc_n^2$ diverges, while $\displaystyle\sum_{n\ge0}\frac{c_n^2}{n+1}=\frac4\pi$.
\end{enumerate}
\end{exercise}

\begin{exercise}[$\star\star$]\label{pb:catalanmoments}
\emph{Catalan numbers are moments.} Let $C_n=\frac1{n+1}\binom{2n}n$.
\begin{enumerate}[label=(\alph*)]
\item Show that $\displaystyle\frac1{2\pi}\int_{-2}^2x^{2n}\sqrt{4-x^2}\dd x=C_n$, while the odd moments vanish.
\item Show that $\displaystyle\frac1{2\pi}\int_0^4x^n\sqrt{\frac{4-x}x}\dd x=C_n$.
\item From (b), show that $C_n\sim\frac{4^n}{\sqrt\pi\,n^{3/2}}$.
\end{enumerate}
\end{exercise}

\section{Fractional parts}

Write $x=\lfloor x\rfloor+\{x\}$. The function $\{1/x\}$ oscillates infinitely often near $0$, yet it is integrable, and its integral
is the first appearance of $1-\gamma$ in this chapter. The same constant governs the remainders of the division of $x$ by
$n=1,2,3,\dots$, and, by the prime number theorem, the remainders of $x$ modulo the primes.

\begin{exercise}[$\star\star$]\label{pb:fracint}
\emph{Integrals of $\{1/x\}$.}
\begin{enumerate}[label=(\alph*)]
\item Show that $\displaystyle\int_0^1\Bigl\{\frac1x\Bigr\}\dd x=1-\gamma$.
\item Show that $\displaystyle\int_0^1\Bigl\{\frac1x\Bigr\}^2\dd x=\log(2\pi)-\gamma-1$.
\item Show that $\displaystyle\int_0^1\frac{\mathrm dx}{\lfloor1/x\rfloor}=\zeta(2)-1$.
\end{enumerate}
\end{exercise}

\begin{exercise}[$\star\star$]\label{pb:fracmean}
\emph{Remainders on average.}
\begin{enumerate}[label=(\alph*)]
\item Show that $\displaystyle\sum_{n\le x}\Bigl\{\frac xn\Bigr\}=(1-\gamma)x+O(\sqrt x)$ for $x\ge1$.
\item For an integer $N\ge2$, let $N\bmod n$ be the remainder of the division of $N$ by $n$. Show that
\[
\sum_{n=1}^N(N\bmod n)=\Bigl(1-\frac{\pi^2}{12}\Bigr)N^2+O(N\log N).
\]
\item Explain both constants as integrals: $1-\gamma=\int_0^1\{\frac1t\}\dd t$ and $2-\frac{\pi^2}6=\int_0^12t\{\frac1t\}\dd t$.
\end{enumerate}
\end{exercise}

\begin{exercise}[$\star\star$]\label{pb:quotientlaw}
\emph{The law of $\{1/U\}$.} Let $U$ be uniform on $]0,1[$.
\begin{enumerate}[label=(\alph*)]
\item Show that $\{1/U\}$ has the density $g(u)=\sum_{k\ge1}\frac1{(k+u)^2}$ on $[0,1]$, which decreases from $\frac{\pi^2}6$ to $\frac{\pi^2}6-1$.
\item Show that $P\bigl(\lfloor1/U\rfloor\text{ is odd}\bigr)=\log2$ and $P\bigl(\{1/U\}<\frac12\bigr)=2-2\log2$.
\item Show that the proportion of $n\le x$ for which $\lfloor x/n\rfloor$ is odd tends to $\log2$ as $x\to\infty$.
\end{enumerate}
\end{exercise}

\begin{figure}[ht]
\centering
\fig{ch09-quotient}
\caption{The density of $\{1/U\}$ (Problem~\ref{pb:quotientlaw}). Small fractional parts are more likely: the shaded area, the chance
that $\{1/U\}<\frac12$, is $2-2\log2=0.614$, and the mean is $1-\gamma=0.423$ instead of $\frac12$. The same law describes the
remainders $\{x/n\}$ for $n\le x$ (Problem~\ref{pb:fracmean}) and, by the prime number theorem, $\{x/p\}$ for primes $p\le x$
(Problem~\ref{pb:primefrac}).}
\label{fig:quotient}
\end{figure}

\begin{exercise}[$\star\star\star$]\label{pb:primefrac}
\emph{Remainders modulo primes (de la Vallée Poussin, 1898).} Show that
\[
\frac1{\pi(x)}\sum_{p\le x}\Bigl\{\frac xp\Bigr\}\longrightarrow1-\gamma\qquad(x\to\infty).
\]
\end{exercise}

\section{Uniform random numbers}

A point chosen uniformly in a cube turns volumes into probabilities, and the volumes of simple polytopes are rational numbers
with factorials in them. In 2001 David and Jonathan Borwein published a family of integrals that seem to follow an obvious pattern
and then, at the eighth step, miss it by about $2\cdot10^{-11}$; the explanation (Problem~\ref{pb:borwein}) is a sum of uniform
random numbers that, for the first time, is able to exceed $1$.

\begin{exercise}[$\star$]\label{pb:passone}
\emph{Waiting to pass $1$.} Let $U_1,U_2,\dots$ be independent and uniform on $[0,1]$.
\begin{enumerate}[label=(\alph*)]
\item Let $N$ be the least $n$ with $U_1+\cdots+U_n>1$. Show that $P(N>n)=\frac1{n!}$, $\mathbb E[N]=e$ and $\operatorname{Var}N=3e-e^2$.
\item Let $L$ be the least $n$ with $U_{n+1}<U_n$, the length of the first increasing run. Show that $L+1$ has the same law as $N$.
\item Find a volume-preserving linear bijection that explains (b) without computation.
\end{enumerate}
\end{exercise}

\begin{exercise}[$\star\star\star$]\label{pb:irwinhall}
\emph{Sums of uniforms and Eulerian numbers.} Let $S_n=U_1+\cdots+U_n$ with independent uniform $U_i$ on $[0,1[$, and let $A(n,k)$
be the number of permutations of $[n]$ with exactly $k$ descents (positions $i$ with $\sigma(i)>\sigma(i+1)$).
\begin{enumerate}[label=(\alph*)]
\item Show that $\displaystyle P(S_n\le x)=\frac1{n!}\sum_{j\ge0}(-1)^j\binom nj(x-j)_+^n$, where $t_+=\max(t,0)$.
\item Let $V_0=0$ and $V_i=\{S_i\}$. Show that $V_1,\dots,V_n$ are independent and uniform on $[0,1[$, and that $\lfloor S_n\rfloor$ is the
number of $i\in\{1,\dots,n-1\}$ with $V_{i+1}<V_i$.
\item Deduce that $P\bigl(\lfloor S_n\rfloor=k\bigr)=\frac{A(n,k)}{n!}$ and $\displaystyle A(n,k)=\sum_{j=0}^k(-1)^j\binom{n+1}j(k+1-j)^n$.
\end{enumerate}
\end{exercise}

\begin{figure}[ht]
\centering
\fig{ch09-irwinhall}
\caption{The densities of $U_1+U_2$, $U_1+U_2+U_3$ and $U_1+\cdots+U_4$. The areas under the last one over $[k,k+1[$ are
$\frac1{24},\frac{11}{24},\frac{11}{24},\frac1{24}$: the Eulerian numbers $1,11,11,1$ count the permutations of $4$ elements with $0,1,2,3$
descents (Problem~\ref{pb:irwinhall}).}
\label{fig:irwinhall}
\end{figure}

\begin{exercise}[$\star\star$]\label{pb:sincpower}
\emph{Powers of $\frac{\sin t}t$.}
\begin{enumerate}[label=(\alph*)]
\item For $n\ge2$, show that
\[
\int_{-\infty}^\infty\Bigl(\frac{\sin t}t\Bigr)^n\dd t=\frac\pi{2^{n-1}(n-1)!}\sum_{0\le k<n/2}(-1)^k\binom nk(n-2k)^{n-1},
\]
and check the values $\pi,\frac{3\pi}4,\frac{2\pi}3$ for $n=2,3,4$.
\item Show that this integral is equivalent to $\sqrt{\frac{6\pi}n}$ as $n\to\infty$.
\end{enumerate}
\end{exercise}

\begin{exercise}[$\star\star\star$]\label{pb:borwein}
\emph{The Borwein integrals.} Write $\operatorname{sinc}x=\frac{\sin x}x$. Let $a_0,a_1,\dots,a_m>0$ with $m\ge1$.
\begin{enumerate}[label=(\alph*)]
\item Let $V_1,\dots,V_m$ be independent and uniform on $[-1,1]$, and $Y=\sum_{k=1}^ma_kV_k$. Show that
\[
\int_0^\infty\prod_{k=0}^m\operatorname{sinc}(a_kx)\dd x=\frac\pi{2a_0}\,P\bigl(|Y|<a_0\bigr).
\]
\item Deduce that $\displaystyle\int_0^\infty\prod_{k=0}^m\operatorname{sinc}\Bigl(\frac x{2k+1}\Bigr)\dd x=\frac\pi2$ for $m=0,1,\dots,6$, but not for $m=7$.
\item For $m=7$, let $\delta=\frac13+\frac15+\cdots+\frac1{15}-1$. Show that the integral equals
\[
\frac\pi2-\frac{\pi\,\delta^7\cdot3\cdot5\cdot7\cdots15}{2^7\cdot7!}=\frac\pi2-2.31\ldots\cdot10^{-11}.
\]
\end{enumerate}
\end{exercise}

\begin{exercise}[$\star$]\label{pb:stick}
\emph{Breaking a stick.} A stick of length $1$ is broken at $n-1$ independent uniform points, $n\ge3$.
\begin{enumerate}[label=(\alph*)]
\item Show that every piece is shorter than $\frac12$ with probability $1-\frac n{2^{n-1}}$. (This is the probability that the pieces
are the sides of an $n$-gon: $\frac14$ for a triangle.)
\item A point $T$ is chosen uniformly on the stick, independently. Show that the piece containing $T$ has expected length $\frac2{n+1}$,
nearly twice the average length $\frac1n$.
\end{enumerate}
\end{exercise}

\begin{exercise}[$\star\star$]\label{pb:spacings}
\emph{The longest and the shortest piece.} Let $L_1,\dots,L_n$ be the lengths of the pieces of Problem~\ref{pb:stick}, from left to right,
and $E_1,\dots,E_n$ independent exponential variables of parameter $1$, with sum $S$.
\begin{enumerate}[label=(\alph*)]
\item Show that $(L_1,\dots,L_n)$ has the same law as $(\frac{E_1}S,\dots,\frac{E_n}S)$, and that this vector is independent of $S$.
\item Deduce that $\mathbb E[\max_iL_i]=\frac{H_n}n$ and $\mathbb E[\min_iL_i]=\frac1{n^2}$, where $H_n=1+\frac12+\cdots+\frac1n$.
\end{enumerate}
\end{exercise}

\begin{exercise}[$\star$]\label{pb:sophomore}
\emph{The sophomore's dream (Johann Bernoulli, 1697).}
\begin{enumerate}[label=(\alph*)]
\item Show that $\displaystyle\int_0^1x^{-x}\dd x=\sum_{n\ge1}\frac1{n^n}$ and $\displaystyle\int_0^1x^x\dd x=\sum_{n\ge1}\frac{(-1)^{n+1}}{n^n}$.
\item How many terms give the first integral to $10$ decimal places?
\end{enumerate}
\end{exercise}

\section{Hints}

\begin{hint}{pb:logtwo}
(a) $\frac{\mathrm d^n}{\mathrm dx^n}x^{n+k}=\frac{(n+k)!}{k!}x^k$, and integrate by parts $n$ times. (b) $\int_0^1\frac{x^k}{1+x}\dd x$ is $\pm\log2$ plus a rational with
denominator dividing $d_k$. (c) Maximise $\frac{x(1-x)}{1+x}$; bound $d_n$ with Chebyshev's $\prod_{p\le x}p<4^x$.
\end{hint}

\begin{hint}{pb:hermite}
Imitate Niven's proof for $\pi^2$ with $f(x)=\frac{x^n(1-x)^n}{n!}$ and $F=\sum_j(-1)^jk^{2n-j}f^{(j)}$: compute $\frac{\mathrm d}{\mathrm dx}\bigl(e^{kx}F(x)\bigr)$.
\end{hint}

\begin{hint}{pb:doubleint}
Expand $\frac1{1-xy}=\sum_k(xy)^k$ and integrate term by term; for (b) differentiate $\int\!\!\int x^{r+t}y^{s+t}/(1-xy)$ in $t$, or use
$\int_0^1x^k(-\log x)\dd x=\frac1{(k+1)^2}$.
\end{hint}

\begin{hint}{pb:beukerstwo}
(a) Problem~\ref{pb:doubleint}(a). (b) For fixed $y$, $\frac{\partial^n}{\partial x^n}\frac1{1-xy}=\frac{n!\,y^n}{(1-xy)^{n+1}}$. (c) $(1-x)(1-y)\le(1-\sqrt{xy})^2$.
(d) $d_n\le3^n$ for large $n$.
\end{hint}

\begin{hint}{pb:beukersthree}
(c) After the substitution, $\frac{(xyz)^n\dd z}{(1-(1-xy)z)^{n+1}}$ becomes $\frac{(1-w)^n\dd w}{1-(1-xy)w}$. (d) First maximise in $w$ for fixed $xy$; the
answer is $(1+\sqrt{xy})^{-2}$. (e) $27(\sqrt2-1)^4<1$.
\end{hint}

\begin{hint}{pb:sondow}
$\frac1{-\log u}=\int_0^\infty u^t\dd t$ for $0<u<1$. You should meet $\sum_{m\ge1}\bigl(\frac1m-\log(1+\frac1m)\bigr)$, and for (b) Wallis's product.
\end{hint}

\begin{hint}{pb:dalzell}
(a) Divide $x^4(1-x)^4$ by $1+x^2$. (b) $x^4(1-x)^4\equiv-4$ modulo $1+x^2$: evaluate at $x=i$.
\end{hint}

\begin{hint}{pb:centralsum}
(a) The Beta integral. (b) $\frac1{\binom{2n}n}=(2n+1)\int_0^1(t-t^2)^n\dd t$ and $\sum_n(2n+1)u^n=\frac{1+u}{(1-u)^2}$.
\end{hint}

\begin{hint}{pb:arcsin}
(b) Find the recurrence for the Taylor coefficients, and for uniqueness note that $(\sqrt{1-x^2}\,D')'=0$ for the difference $D$.
(c) $x=\frac12$ and $x=1$.
\end{hint}

\begin{hint}{pb:lehmer}
(a) $\sum c_nx^{2n}=\frac1{\sqrt{1-x^2}}$. (b) $c_{n+1}=\frac{2n+1}{2n+2}c_n$; look for $b_n=\lambda nc_n^2$ with $b_n-b_{n+1}=\frac{c_n^2}{n+1}$.
\end{hint}

\begin{hint}{pb:catalanmoments}
(a) $x=2\sin\theta$ and the Beta integral. (b) $x=y^2$. (c) Laplace's method near $x=4$.
\end{hint}

\begin{hint}{pb:fracint}
Substitute $x=\frac1t$ and cut $[1,\infty[$ into the intervals $[k,k+1[$. For (b) you will need Stirling's formula.
\end{hint}

\begin{hint}{pb:fracmean}
(a) $\sum_{n\le x}\lfloor\frac xn\rfloor=\sum_{m\le x}\tau(m)$. (b) $\sum_{n\le N}n\lfloor\frac Nn\rfloor=\sum_{m\le N}\sigma(m)$.
\end{hint}

\begin{hint}{pb:quotientlaw}
(a) $\{1/U\}\le u$ and $\lfloor 1/U\rfloor=k$ mean $\frac1{k+u}\le U\le\frac1k$. (c) Count the $n$ with $\lfloor x/n\rfloor=k$, for $k\le K$.
\end{hint}

\begin{hint}{pb:primefrac}
The function $\{1/t\}$ is Riemann integrable on $[\frac1K,1]$; sandwich it between step functions and count primes in each interval
$]xt_{j-1},xt_j]$ by the prime number theorem. The primes below $x/K$ are few.
\end{hint}

\begin{hint}{pb:passone}
(a) $P(N>n)=P(U_1+\cdots+U_n\le1)$, and $\mathbb E[N^2]=\sum_{n\ge0}(2n+1)P(N>n)$. (c) Partial sums.
\end{hint}

\begin{hint}{pb:irwinhall}
(a) Inclusion--exclusion on the set of coordinates that exceed $1$. (b) Given $V_1,\dots,V_{i-1}$, the map $u\mapsto\{V_{i-1}+u\}$ is a
rotation of $[0,1[$. (c) The relative order of $V_1,\dots,V_n$ is a uniform random permutation.
\end{hint}

\begin{hint}{pb:sincpower}
(a) $\frac{\sin t}t$ is the characteristic function of the uniform law on $[-1,1]$; invert at $0$. (b) Laplace's method on $[-\pi,\pi]$,
where $\log\frac{\sin t}t=-\frac{t^2}6+O(t^4)$.
\end{hint}

\begin{hint}{pb:borwein}
(a) $\prod_{k\ge1}\operatorname{sinc}(a_kx)=\mathbb E[e^{ixY}]$, and $\int_{-R}^R\operatorname{sinc}(a_0x)\cos(xy)\dd x$ is a difference of two values of
$\int_0^z\frac{\sin u}u\dd u$. (c) $Y>1$ is a small corner of the cube, a simplex.
\end{hint}

\begin{hint}{pb:stick}
(a) Put the stick on a circle, with one more uniform point as the place of the break. (b) $\mathbb E\bigl[\sum_iL_i^2\bigr]$.
\end{hint}

\begin{hint}{pb:spacings}
(a) Change variables $(e_1,\dots,e_n)\mapsto(s,\frac{e_1}s,\dots,\frac{e_{n-1}}s)$; the Jacobian is $s^{n-1}$. (b) $\mathbb E[\max E_i]=H_n$ and
$\max E_i=S\cdot\max L_i$ in law.
\end{hint}

\begin{hint}{pb:sophomore}
$x^{-x}=e^{-x\log x}$; expand and use $\int_0^1x^k(-\log x)^k\dd x=\frac{k!}{(k+1)^{k+1}}$.
\end{hint}

\section{Solutions}

\begin{solution}{pb:logtwo}
\idea A polynomial with integer coefficients, divided by $1+x$, integrates to a rational number plus a multiple of $\log2$, and
the denominators of the rationals are controlled by $d_n$. Rodrigues's formula $P_n=\frac1{n!}D^n(x^n(1-x)^n)$ makes the integral small:
the $n$ derivatives can be thrown onto $\frac1{1+x}$.

\writeup (a) $x^n(1-x)^n=\sum_k(-1)^k\binom nkx^{n+k}$, and $\frac1{n!}D^nx^{n+k}=\binom{n+k}nx^k$, so $P_n(x)=\sum_k(-1)^k\binom nk\binom{n+k}kx^k$ has integer
coefficients. The polynomial $x^n(1-x)^n$ has zeros of order $n$ at $0$ and $1$, so its derivatives of order $<n$ vanish there and $n$
integrations by parts give $\int_0^1P_nf=\frac{(-1)^n}{n!}\int_0^1x^n(1-x)^nf^{(n)}$ for every smooth $f$. With $f(x)=\frac1{1+x}$,
$f^{(n)}(x)=\frac{(-1)^nn!}{(1+x)^{n+1}}$, which gives the second formula.

(b) Dividing $x^k$ by $1+x$,
\[
\frac{x^k}{1+x}=\sum_{j=1}^k(-1)^{j-1}x^{k-j}+\frac{(-1)^k}{1+x},
\]
so
\[
\int_0^1\frac{x^k}{1+x}\dd x=(-1)^k\log2+\sum_{j=1}^k\frac{(-1)^{j-1}}{k-j+1},
\]
and $d_k$, hence $d_n$ for $k\le n$, times the last sum is an integer. Summing against the integer coefficients of $P_n$ gives
$d_nI_n=a_n+b_n\log2$ with $b_n=d_nP_n(-1)$ and $a_n\in\Z$.

(c) Let $g(x)=\frac{x(1-x)}{1+x}$. Then $g'(x)=\frac{1-2x-x^2}{(1+x)^2}$ vanishes on $[0,1]$ only at $x=\sqrt2-1$, where $g=(\sqrt2-1)^2=3-2\sqrt2$, its
maximum. So $0<I_n=\int_0^1\frac{g(x)^n}{1+x}\dd x\le(3-2\sqrt2)^n\log2$. Next, $d_n=\prod_{p\le n}p^{\lfloor\log n/\log p\rfloor}$; the primes $p>\sqrt n$ appear
to the first power, and the others contribute at most $n$ each, so by Tome~1, Theorem~2.4(i), $d_n\le4^nn^{\sqrt n}$. Suppose
$\log2=\frac pq$. Then $q(a_n+b_n\log2)=qa_n+pb_n$ is a non-zero integer (as $I_n>0$), and
\[
1\le|qd_nI_n|\le q\,n^{\sqrt n}\bigl(4(3-2\sqrt2)\bigr)^n\log2.
\]
Since $4(3-2\sqrt2)=0.686\ldots<1$ and $n^{\sqrt n}=e^{\sqrt n\log n}$ grows more slowly than any exponential, the right side tends to $0$:
a contradiction.

\tools{Tome~1, Theorem~2.4 (Chebyshev's bound); Tome~0, Chapter~6 (integration by parts); Tome~4, Chapter~4.}
\leads The same integrals give an \emph{irrationality measure}: $|\log2-\frac pq|>q^{-\mu}$ for large $q$. With $d_n=e^{n(1+o(1))}$ and
$|P_n(-1)|=(3+2\sqrt2)^{n(1+o(1))}$, the argument gives $\mu=1+\frac{1+\log(3+2\sqrt2)}{\log(3+2\sqrt2)-1}=4.62\ldots$; refinements of the integrals (Rukhadze, Hata, Marcovecchio) brought $\mu$ down to $3.57\ldots$ (R.~Marcovecchio,
2009). For comparison, almost every real number has $\mu=2$.
\end{solution}

\begin{solution}{pb:hermite}
\idea A polynomial $f$ whose derivatives at $0$ and $1$ are all integers, and an antiderivative of $e^{kx}f(x)$ of the shape $e^{kx}F(x)$,
turn $\int_0^1e^{kx}f$ into an integer combination of $1$ and $e^k$; and the factor $\frac1{n!}$ in $f$ makes the integral tiny.

\writeup (a) Let $f(x)=\frac{x^n(1-x)^n}{n!}=\frac1{n!}\sum_{j=n}^{2n}c_jx^j$ with $c_j\in\Z$. Then $f^{(j)}(0)=\frac{j!}{n!}c_j\in\Z$ for $n\le j\le2n$, and $0$ for
other $j$; since $f(1-x)=f(x)$, also $f^{(j)}(1)\in\Z$. Let $F=\sum_{j=0}^{2n}(-1)^jk^{2n-j}f^{(j)}$. Then
\[
kF+F'=\sum_{j=0}^{2n}(-1)^jk^{2n+1-j}f^{(j)}+\sum_{j=0}^{2n}(-1)^jk^{2n-j}f^{(j+1)}=k^{2n+1}f,
\]
because the sums telescope and $f^{(2n+1)}=0$. So $\frac{\mathrm d}{\mathrm dx}\bigl(e^{kx}F(x)\bigr)=k^{2n+1}e^{kx}f(x)$, and
\[
k^{2n+1}\int_0^1e^{kx}f(x)\dd x=e^kF(1)-F(0).
\]
If $e^k=\frac ab$ with $a,b\ge1$ integers, then
\[
M_n=bk^{2n+1}\int_0^1e^{kx}f(x)\dd x=aF(1)-bF(0)
\]
is an integer. It is positive, since
$f>0$ on $]0,1[$. But $0<f\le\frac{4^{-n}}{n!}$, so $M_n\le\frac{bk e^k(k^2/4)^n}{n!}\to0$. For large $n$, $0<M_n<1$: a contradiction.

(b) Let $r=\frac pq$ with $p\ne0$, $q\ge1$. If $e^r$ were rational, so would be $(e^r)^q=e^p$ and $e^{-p}=\frac1{e^p}$, one of which is $e^k$
with $k=|p|\ge1$; this contradicts (a). If $\log r=\frac pq$ for a rational $r>0$, $r\ne1$, then $p\ne0$ and $e^{p/q}=r$ would be
rational.

\tools{Tome~4, Theorem~3.8 (the same method for $\pi^2$); Tome~0, Exercise~6.4 ($e$ is irrational).}
\leads Hermite (1873) pushed this idea, with several functions $e^{kx}$ at once, to prove that $e$ is transcendental: it is the root of
no non-zero polynomial with integer coefficients. Lindemann (1882) proved that $\pi$ is transcendental, which shows that the circle
cannot be squared with ruler and compass.
\end{solution}

\begin{solution}{pb:doubleint}
\idea Expand $\frac1{1-xy}$ as a geometric series; each term is a product of two one-variable integrals, and the sums telescope.

\writeup (a) The terms of $\sum_k(xy)^kx^ry^s$ are non-negative, so we may integrate term by term (Tome~4, Chapter~4):
\[
\int_0^1\!\!\int_0^1\frac{x^ry^s}{1-xy}\dd x\dd y=\sum_{k\ge0}\frac1{(r+k+1)(s+k+1)}.
\]
For $r>s$,
\[
\frac1{(r+k+1)(s+k+1)}=\frac1{r-s}\Bigl(\frac1{s+k+1}-\frac1{r+k+1}\Bigr),
\]
and the sum telescopes to $\frac1{r-s}\sum_{j=s+1}^r\frac1j$. For
$r=s$ the sum is $\sum_{j>r}\frac1{j^2}=\zeta(2)-\sum_{j\le r}\frac1{j^2}$.

(b) We have $-\log(xy)=-\log x-\log y$, and integration by parts gives
\[
\int_0^1x^a(-\log x)\dd x=\frac1{(a+1)^2}.
\]
So the term $(xy)^kx^ry^s$ contributes $\frac1{A^2B}+\frac1{AB^2}$, where $A=r+k+1$ and $B=s+k+1$. For $r>s$ this equals
$\frac1{r-s}\bigl(\frac1{B^2}-\frac1{A^2}\bigr)$, because $\frac1{B^2}-\frac1{A^2}=\frac{(A-B)(A+B)}{A^2B^2}$ and $A-B=r-s$. The sum telescopes again. For $r=s$ the term is
$\frac2{(r+k+1)^3}$, and the sum is $2\sum_{j>r}\frac1{j^3}$.

(c) $\frac1{1+xy}=\sum_k(-xy)^k$, and the partial sums lie between $0$ and $1$, so term-by-term integration is legitimate (dominated
convergence, Tome~4, Chapter~4). The result is $\sum_{k\ge0}\frac{(-1)^k}{(k+1)^2}=\zeta(2)-2\cdot\frac{\zeta(2)}4=\frac{\pi^2}{12}$, by Tome~4, Theorem~3.7.

\tools{Tome~4, Chapter~4 (term-by-term integration) and Theorem~3.7 ($\zeta(2)=\frac{\pi^2}6$).}
\leads The denominators are the point: for $r,s\le n$ the rational parts in (a) become integers when multiplied by $d_n^2$, and those
in (b) when multiplied by $d_n^3$. This is what Problems~\ref{pb:beukerstwo} and~\ref{pb:beukersthree} use.
\end{solution}

\begin{solution}{pb:beukerstwo}
\idea Problem~\ref{pb:doubleint} shows that any polynomial in $x$ and $y$ of degree at most $n$ in each, integrated against $\frac1{1-xy}$,
gives a rational number with denominator $d_n^2$ plus a rational multiple of $\zeta(2)$. Choosing the polynomial $(1-y)^nP_n(x)$ makes
the integral small, by the same trick as for $\log2$.

\writeup (a) Expand $(1-y)^nP_n(x)=\sum_{r,s\le n}p_rq_sx^ry^s$ with integers $p_r,q_s$. By Problem~\ref{pb:doubleint}(a), the term with $r\ne s$
contributes $\frac{p_rq_s}{|r-s|}\sum_{j=\min+1}^{\max}\frac1j$; here $|r-s|$ divides $d_n$ and $d_n\sum\frac1j\in\Z$, so $d_n^2$ times it is an
integer. The term with $r=s$ contributes $p_rq_r\bigl(\zeta(2)-\sum_{j\le r}\frac1{j^2}\bigr)$, and $d_n^2\sum_{j\le r}\frac1{j^2}\in\Z$. So
$d_n^2J_n=a_n+b_n\zeta(2)$ with $b_n=d_n^2\sum_rp_rq_r$ and $a_n\in\Z$.

(b) For fixed $y\in[0,1[$, the function $f(x)=\frac1{1-xy}$ is smooth on $[0,1]$ with $f^{(n)}(x)=\frac{n!\,y^n}{(1-xy)^{n+1}}$. As in
Problem~\ref{pb:logtwo}(a),
\[
\int_0^1\frac{P_n(x)}{1-xy}\dd x=\frac{(-1)^n}{n!}\int_0^1x^n(1-x)^n\frac{n!\,y^n}{(1-xy)^{n+1}}\dd x.
\]
Multiply by $(1-y)^n$ and integrate over $y\in[0,1[$.

(c) Let $t=\sqrt{xy}$. By AM--GM, $x+y\ge2t$, so $(1-x)(1-y)=1-(x+y)+xy\le(1-t)^2$. Hence
\[
\frac{x(1-x)y(1-y)}{1-xy}\le\frac{t^2(1-t)^2}{1-t^2}=\frac{t^2(1-t)}{1+t}=h(t).
\]
On $]0,1[$, $(\log h)'=\frac2t-\frac1{1-t}-\frac1{1+t}$ vanishes when $2(1-t^2)=2t$, that is $t^2+t-1=0$, $t=\tau=\frac{\sqrt5-1}2$, and $h$
is maximal there. Since $\tau^2=1-\tau$ and $\tau(1+\tau)=1$, $h(\tau)=\tau^2\cdot\tau^2\cdot\tau=\tau^5$. Therefore, by (b),
\[
0<|J_n|\le\tau^{5n}\int_0^1\!\!\int_0^1\frac{\mathrm dx\dd y}{1-xy}=\zeta(2)\tau^{5n}.
\]

(d) By the prime number theorem (Tome~5, Theorem~5.5), $\log d_n=\psi(n)\sim n$, so $d_n\le3^n$ for $n$ large. If $\zeta(2)=\frac pq$, then
$q d_n^2J_n=qa_n+pb_n$ is a non-zero integer with $|qd_n^2J_n|\le q\zeta(2)(9\tau^5)^n$, and $9\tau^5=0.81\ldots<1$: a contradiction for large $n$.

\tools{Problems~\ref{pb:logtwo} and~\ref{pb:doubleint}; Tome~0, Theorem~4.3 (AM--GM); Tome~5, Theorem~5.5 (prime number theorem).}
\leads The numbers $b_n$ are, up to sign and the factor $d_n^2$, the integers $\sum_k\binom nk^2\binom{n+k}k=1,3,19,147,1251,\dots$ that Apéry also
used for $\zeta(2)$; they satisfy a three-term recurrence, checked numerically. Tome~4 proved the irrationality of
$\pi^2$ by Niven's method, which says nothing about $\zeta(3)$; Beukers's method does.
\end{solution}

\begin{solution}{pb:beukersthree}
\idea Same plan as for $\zeta(2)$, one dimension up. The kernel $\frac{-\log(xy)}{1-xy}$ produces $\zeta(3)$ with denominator $d_n^3$
(Problem~\ref{pb:doubleint}(b)). Writing the logarithm as an integral in a third variable gives room for a change of variables that
lets us integrate by parts in both $x$ and $y$.

\writeup (a) As in Problem~\ref{pb:beukerstwo}(a), using Problem~\ref{pb:doubleint}(b): for $r\ne s$, $|r-s|$ divides $d_n$ and $d_n^2\sum\frac1{j^2}\in\Z$; for
$r=s$, $d_n^3\sum_{j\le r}\frac1{j^3}\in\Z$.

(b) $\int_0^1\frac{\mathrm dz}{1-(1-xy)z}=\Bigl[\frac{-\log(1-(1-xy)z)}{1-xy}\Bigr]_0^1=\frac{-\log(xy)}{1-xy}$.

(c) So $K_n=\int\!\!\int\!\!\int\frac{P_n(x)P_n(y)}{1-(1-xy)z}\dd x\dd y\dd z$ (the integrand is positive up to the bounded factor $P_n(x)P_n(y)$, and
$\int\!\!\int\frac{-\log(xy)}{1-xy}=2\zeta(3)$ is finite, so the order of integration may be changed). For fixed $y,z$, $f(x)=\frac1{1-z+xyz}$ has
$f^{(n)}=\frac{(-1)^nn!\,(yz)^n}{(1-(1-xy)z)^{n+1}}$, and $n$ integrations by parts in $x$ give
\[
K_n=\int\!\!\int\!\!\int\frac{(xyz)^n(1-x)^nP_n(y)}{(1-(1-xy)z)^{n+1}}\dd x\dd y\dd z.
\]
Now fix $x,y$ and put $w=\frac{1-z}{1-(1-xy)z}$, a decreasing bijection of $[0,1]$. Then $z=\frac{1-w}{1-(1-xy)w}$,
$1-(1-xy)z=\frac{xy}{1-(1-xy)w}$ and $\mathrm dz=-\frac{xy\dd w}{(1-(1-xy)w)^2}$, so
\[
\frac{(xyz)^n\dd z}{(1-(1-xy)z)^{n+1}}=-\frac{(1-w)^n\dd w}{1-(1-xy)w}
\]
(the powers of $xy$ and of $1-(1-xy)w$ cancel). Hence
\[
K_n=\int\!\!\int\!\!\int\frac{(1-x)^n(1-w)^nP_n(y)}{1-(1-xy)w}\dd x\dd y\dd w.
\] For fixed $x,w$,
$g(y)=\frac1{1-w+xyw}$ has $g^{(n)}=\frac{(-1)^nn!\,(xw)^n}{(1-(1-xy)w)^{n+1}}$, and $n$ integrations by parts in $y$ give the formula.

(d) Fix $x,y\in]0,1]$ and let $\rho=\sqrt{xy}$, $s=1-xy=1-\rho^2$. The function $w\mapsto\frac{w(1-w)}{1-sw}$ has derivative of the sign of
$(1-2w)(1-sw)+sw(1-w)=1-2w+sw^2$, which vanishes at $w=\frac{1-\rho}{s}=\frac1{1+\rho}$; there $1-w=\frac\rho{1+\rho}$, $1-sw=\rho$, and the value
is $\frac1{(1+\rho)^2}$, the maximum. As in Problem~\ref{pb:beukerstwo}(c), $x(1-x)y(1-y)\le\rho^2(1-\rho)^2$, so the whole expression is at most
\[
\Bigl(\frac{\rho(1-\rho)}{1+\rho}\Bigr)^2\le(3-2\sqrt2)^2=(\sqrt2-1)^4,
\]
by the maximum of $\frac{t(1-t)}{1+t}$ found in Problem~\ref{pb:logtwo}(c). By (c) with $n=0$, the integral of $\frac1{1-(1-xy)w}$ over the cube is
$K_0=\int\!\!\int\frac{-\log(xy)}{1-xy}=2\zeta(3)$ (Problem~\ref{pb:doubleint}(b)). So $0<K_n\le2\zeta(3)(\sqrt2-1)^{4n}$.

(e) For large $n$, $d_n\le3^n$, and $27(\sqrt2-1)^4=0.79\ldots<1$. If $\zeta(3)=\frac pq$, the non-zero integers $qd_n^3K_n$ would tend to $0$.

\tools{Problems~\ref{pb:logtwo}, \ref{pb:doubleint} and~\ref{pb:beukerstwo}; Tome~5, Theorem~5.5; Tome~4, Chapter~4 (changing the order of integration).}
\leads The coefficient $b_n$ is $2d_n^3$ times the Apéry number $\sum_k\binom nk^2\binom{n+k}k^2=1,5,73,1445,\dots$, which satisfies
$n^3u_n=(34n^3-51n^2+27n-5)u_{n-1}-(n-1)^3u_{n-2}$; Apéry's original proof ran on this recurrence. The method gives an irrationality
measure $13.41\ldots$ for $\zeta(3)$, later improved to $5.51\ldots$ by Rhin and Viola (2001). For $\zeta(5)$ the analogous integrals do
not decay fast enough against $d_n^5$, and Rivoal's and Zudilin's theorems use cleverer integrals of the same family.
\end{solution}

\begin{solution}{pb:sondow}
\idea The factor $\frac1{-\log(xy)}$ is the integral of $(xy)^t$ over $t\ge0$. With it, everything becomes a sum of rational
functions of $t$, which integrate to logarithms; and $\gamma$ is the limit of $H_N-\log N$, a sum of $\frac1m-\log(1+\frac1m)$.

\writeup (a) For $0<u<1$, $\int_0^\infty u^t\dd t=\frac1{-\log u}$. All terms below are non-negative, so by Tonelli's principle for
non-negative series and integrals (Tome~4, Chapter~4) we may expand and exchange freely:
\[
\int_0^1\!\!\int_0^1\frac{1-x}{(1-xy)(-\log xy)}=\int_0^\infty\sum_{k\ge0}\int_0^1\!\!\int_0^1(1-x)(xy)^{k+t}\dd x\dd y\dd t.
\]
With $a=k+t$, $\int_0^1(1-x)x^a\dd x=\frac1{(a+1)(a+2)}$ and $\int_0^1y^a\dd y=\frac1{a+1}$. So the integral is
$\sum_{k\ge0}\int_{k+1}^\infty\frac{\mathrm du}{u^2(u+1)}$ ($u=a+1$). Since $\frac1{u^2(u+1)}=\frac1{u^2}-\frac1u+\frac1{u+1}$,
\[
\int_m^\infty\frac{\mathrm du}{u^2(u+1)}=\frac1m-\log\Bigl(1+\frac1m\Bigr),
\]
and $\sum_{m=1}^M\bigl(\frac1m-\log(1+\frac1m)\bigr)=H_M-\log(M+1)\to\gamma$ (Tome~4, Corollary~11.6).

(b) With $\frac1{1+xy}=\sum_k(-xy)^k$ the same computation gives $\sum_{m\ge1}(-1)^{m-1}\bigl(\frac1m-\log(1+\frac1m)\bigr)$; the exchange is justified
because the terms decrease in absolute value and the alternating partial sums are dominated. Now $\sum(-1)^{m-1}\frac1m=\log2$, and
\[
\sum_{m=1}^{2M}(-1)^{m-1}\log\frac{m+1}m=\log\prod_{j=1}^M\frac{2j}{2j-1}\cdot\frac{2j}{2j+1}\longrightarrow\log\frac\pi2
\]
by Wallis's product (Tome~4, Corollary~3.6). The result is $\log2-\log\frac\pi2=\log\frac4\pi$.

\tools{Tome~4, Corollaries~11.6 (Euler's constant) and~3.6 (Wallis), Chapter~4.}
\leads Sondow found these formulas while looking for a Beukers-type proof that $\gamma$ is irrational. Replacing $1-x$ by suitable
polynomials gives integrals equal to $a_n+b_n\gamma+c_n\log(\text{integer})$, and Sondow (2003) showed that $\gamma$ is irrational if and
only if a certain explicit sequence of such numbers is not eventually an integer. Nobody has been able to decide it.
\end{solution}

\begin{solution}{pb:dalzell}
\idea Long division of the numerator by $1+x^2$ leaves a polynomial, whose integral is rational, and a constant remainder, whose
integral against $\frac1{1+x^2}$ is a multiple of $\frac\pi4$.

\writeup (a) One checks that $x^4(1-x)^4=(x^6-4x^5+5x^4-4x^2+4)(1+x^2)-4$. So the integral is
$\frac17-\frac46+1-\frac43+4-4\cdot\frac\pi4=\frac{22}7-\pi$. Since $1\le1+x^2\le2$ on $[0,1]$ and $\int_0^1x^4(1-x)^4\dd x=\frac{4!\,4!}{9!}=\frac1{630}$ (Beta
integral), the integral lies strictly between $\frac1{1260}$ and $\frac1{630}$.

(b) At $x=i$, $x^4(1-x)^4=(1-i)^4=(-2i)^2=-4$, and the same at $x=-i$. So the real polynomial $x^4(1-x)^4+4$ vanishes at $\pm i$ and is
divisible by $1+x^2$. Writing $A=x^4(1-x)^4$, $A^n-(-4)^n=(A+4)(A^{n-1}+\cdots+(-4)^{n-1})$ is divisible by $1+x^2$ too: $A^n=(-4)^n+(1+x^2)Q_n$ with
$Q_n\in\Z[x]$. So the integral equals $r_n+(-4)^n\frac\pi4$ with $r_n=\int_0^1Q_n\in\Q$. It is positive and at most $\int_0^1(x(1-x))^{4n}\le4^{-4n}$.
Multiplying by $4(-4)^{-n}$, $|\pi+4(-4)^{-n}r_n|\le4\cdot4^{-n}\cdot4^{-4n}=4^{1-5n}$, with strict inequality.

\tools{Tome~4, Chapter~5 (Beta integral); Tome~0, Chapter~3 (division of polynomials).}
\leads D.~P.~Dalzell published (a) in 1944. Stephen Lucas (2009) studied the whole family; the approximations in (b) converge
geometrically, but the denominators of $r_n$ grow too fast for an irrationality proof. Mahler (1953) showed by much finer integrals
that $|\pi-\frac pq|>q^{-42}$ for large $q$; the best exponent known today is $7.10\ldots$ (Zeilberger and Zudilin, 2020).
\end{solution}

\begin{solution}{pb:centralsum}
\idea $\frac{(n-1)!\,n!}{(2n)!}$ is a Beta integral, and summing the integrands is a geometric series in $t(1-t)$, which is at most
$\frac14$. What remains is a rational function whose denominator $1-t+t^2$ has roots at primitive sixth roots of unity: hence $\pi\sqrt3$.

\writeup (a) By Tome~4, Theorem~5.11, $\int_0^1t^{n-1}(1-t)^n\dd t=B(n,n+1)=\frac{(n-1)!\,n!}{(2n)!}=\frac1{n\binom{2n}n}$. Summing the non-negative
integrands, $\sum_{n\ge1}t^{n-1}(1-t)^n=\frac{1-t}{1-t+t^2}$. By the symmetry $t\mapsto1-t$, which fixes $1-t+t^2$, its integral is half of
\[
\int_0^1\frac{\mathrm dt}{1-t+t^2}=\frac2{\sqrt3}\Bigl[\arctan\frac{2t-1}{\sqrt3}\Bigr]_0^1=\frac2{\sqrt3}\cdot\frac\pi3=\frac{2\pi}{3\sqrt3}.
\]
So the sum is $\frac\pi{3\sqrt3}=0.6046\ldots$

(b) $\frac1{\binom{2n}n}=(2n+1)B(n+1,n+1)=(2n+1)\int_0^1u^n\dd t$ with $u=t-t^2$, and $\sum_{n\ge0}(2n+1)u^n=\frac{1+u}{(1-u)^2}$ (differentiate the geometric
series). Let $q=1-t+t^2=1-u$. Then $\frac{1+u}{(1-u)^2}=\frac{2-q}{q^2}=\frac2{q^2}-\frac1q$. With $t-\frac12=\frac{\sqrt3}2\tan\theta$, $q=\frac34\sec^2\theta$ and
$\theta\in[-\frac\pi6,\frac\pi6]$:
\[
\int_0^1\frac{\mathrm dt}{q^2}=\frac{8\sqrt3}9\int_{-\pi/6}^{\pi/6}\cos^2\theta\dd\theta=\frac{8\sqrt3}9\Bigl(\frac\pi6+\frac{\sqrt3}4\Bigr)=\frac{4\pi\sqrt3}{27}+\frac23.
\]
So $\sum_{n\ge0}\frac1{\binom{2n}n}=\frac{8\pi\sqrt3}{27}+\frac43-\frac{2\pi}{3\sqrt3}=\frac43+\frac{2\pi\sqrt3}{27}=1.7364\ldots$

\tools{Tome~4, Theorem~5.11 (Beta and Gamma); Tome~0, Chapter~6 (integrals of rational functions).}
\leads D.~H.~Lehmer's 1985 article collects dozens of such sums. The appearance of $\sqrt3$ is no accident: these series are values
of $L$-functions of characters modulo $3$, and the same technique with $\frac1{n^3\binom{2n}n}$ leads to Apéry's series for $\zeta(3)$.
\end{solution}

\begin{solution}{pb:arcsin}
\idea A function satisfying a linear differential equation with polynomial coefficients has Taylor coefficients satisfying a
recurrence. Here the recurrence is a ratio of quadratics, exactly the ratio of consecutive terms of $\frac{4^n}{n^2\binom{2n}n}$.

\writeup (a) $y'=\frac{2\arcsin x}{\sqrt{1-x^2}}$, so $\sqrt{1-x^2}\,y'=2\arcsin x$; differentiating, $\sqrt{1-x^2}\,y''-\frac x{\sqrt{1-x^2}}y'=\frac2{\sqrt{1-x^2}}$, and
multiply by $\sqrt{1-x^2}$.

(b) Let $S(x)=\sum_{n\ge1}c_{2n}x^{2n}$ with $c_{2n}=\frac{4^n}{2n^2\binom{2n}n}$. Since $\frac{c_{2n+2}}{c_{2n}}=\frac{(2n)^2}{(2n+1)(2n+2)}\to1$, the radius of
convergence is $1$. Writing $(1-x^2)S''-xS'=\sum_m\bigl((m+2)(m+1)c_{m+2}-m^2c_m\bigr)x^m$ (with $c_m=0$ for odd $m$ and $c_0=0$), the coefficient
of $x^0$ is $2c_2=2$ and the others vanish by the ratio just computed. So $S$ solves the same equation as $y$ with the same initial values.
For $D=y-S$, $(1-x^2)D''-xD'=0$ means $\bigl(\sqrt{1-x^2}\,D'\bigr)'=0$, so $\sqrt{1-x^2}\,D'$ is constant, equal to its value $0$ at $0$; hence $D'=0$
and $D=0$ on $]-1,1[$. At $x=\pm1$: the terms are non-negative and $c_{2n}\sim\frac{\sqrt{\pi}}2n^{-3/2}$, so the series converges at $1$,
and by monotone convergence $S(x)\to S(1)$ as $x\to1^-$, while $\arcsin^2x\to\frac{\pi^2}4$.

(c) At $x=\frac12$, $\arcsin^2\frac12=\frac{\pi^2}{36}=\frac12\sum\frac1{n^2\binom{2n}n}$. At $x=1$, $\frac{\pi^2}4=\frac12\sum\frac{4^n}{n^2\binom{2n}n}$.

\tools{Tome~4, Chapter~2 (power series) and Proposition~11.12 ($\binom{2n}n\sim\frac{4^n}{\sqrt{\pi n}}$); Tome~0, Chapter~6 (derivatives).}
\leads Euler knew (c). The same idea for $\frac1{n^4\binom{2n}n}$ gives Comtet's $\zeta(4)=\frac{36}{17}\sum_{n\ge1}\frac1{n^4\binom{2n}n}$ (1974). Integer-relation
searches by Borwein and Bradley (1997) found no formula of the same simple shape for $\zeta(5)$, which suggests that $\zeta(3)$ is special.
\end{solution}

\begin{solution}{pb:lehmer}
\idea Everything follows from $\sum_nc_nx^n=\frac1{\sqrt{1-x}}$, the binomial series for the exponent $-\frac12$. Part (b) hides a telescoping
sum, found by trying a multiple of $nc_n^2$.

\writeup (a) For $|x|<1$, $\sum_nc_nx^{2n}=(1-x^2)^{-1/2}$; integrating from $0$ to $x$, $\arcsin x=\sum_n\frac{c_n}{2n+1}x^{2n+1}$. The terms are
non-negative and $c_n\le1$, so the series converges at $x=1$ and, by monotone convergence, its sum is $\arcsin1=\frac\pi2$. Next,
$\sum_{n\ge1}c_nx^{n-1}=\frac{(1-x)^{-1/2}-1}x$. Integrating from $0$ to $1$ and substituting $v=\sqrt{1-x}$ ($x=1-v^2$, $\mathrm dx=-2v\dd v$):
\[
\sum_{n\ge1}\frac{c_n}n=\int_0^1\frac{1/v-1}{1-v^2}\,2v\dd v=\int_0^1\frac{2\dd v}{1+v}=2\log2,
\]
again with monotone convergence at the end point.

(b) $c_n\sim\frac1{\sqrt{\pi n}}$ (Tome~4, Proposition~11.12), so $c_n^2\sim\frac1{\pi n}$ and $\sum c_n^2$ diverges like the harmonic series. Let $b_n=-4nc_n^2$. Using
$c_{n+1}=\frac{2n+1}{2n+2}c_n$,
\[
b_n-b_{n+1}=4c_n^2\Bigl(\frac{(2n+1)^2}{4(n+1)}-n\Bigr)=4c_n^2\cdot\frac1{4(n+1)}=\frac{c_n^2}{n+1}.
\]
So $\sum_{n=0}^{N-1}\frac{c_n^2}{n+1}=b_0-b_N=4Nc_N^2\to\frac4\pi$.

\tools{Tome~0, Chapter~6 (Taylor series); Tome~4, Proposition~11.12 (central binomial coefficient).}
\leads $\sum c_n^2x^n=\frac2\pi K(\sqrt x)$, a complete elliptic integral, whose behaviour near $x=1$ explains the divergence (Tome~4, Chapter~12, on the
AGM, meets the same function). Ramanujan's series for $\frac1\pi$, such as $\sum_n\binom{2n}n^3\frac{42n+5}{2^{12n+4}}=\frac1\pi$, belong to this family.
\end{solution}

\begin{solution}{pb:catalanmoments}
\idea Under $x=2\sin\theta$ the moment becomes a Wallis-type integral, which the Beta function evaluates.

\writeup (a) Odd moments vanish by symmetry. With $x=2\sin\theta$, $\sqrt{4-x^2}=2\cos\theta$ and $\mathrm dx=2\cos\theta\dd\theta$:
\[
\frac1{2\pi}\int_{-2}^2x^{2n}\sqrt{4-x^2}\dd x=\frac{4^{n+1}}{\pi}\int_0^{\pi/2}\sin^{2n}\theta\cos^2\theta\dd\theta=\frac{4^{n+1}}{2\pi}B\Bigl(n+\frac12,\frac32\Bigr)
\]
by Tome~4, Exercise~5.3. Now $B(n+\frac12,\frac32)=\frac{\Gamma(n+\frac12)\Gamma(\frac32)}{\Gamma(n+2)}$, with $\Gamma(n+\frac12)=\frac{(2n)!}{4^nn!}\sqrt\pi$ and $\Gamma(\frac32)=\frac{\sqrt\pi}2$, so
the moment is $\frac{4^{n+1}}{2\pi}\cdot\frac{(2n)!\,\pi}{2\cdot4^nn!\,(n+1)!}=\frac{(2n)!}{n!\,(n+1)!}=C_n$.

(b) With $x=y^2$, $\mathrm dx=2y\dd y$ and $\sqrt{\frac{4-x}x}\,\mathrm dx=2\sqrt{4-y^2}\dd y$; so the integral is $\frac1{2\pi}\int_0^22y^{2n}\sqrt{4-y^2}\dd y$, which is (a).

(c) Put $x=4(1-u)$: $C_n=\frac{2\cdot4^n}\pi\int_0^1(1-u)^{n-1/2}\sqrt u\dd u$. As $(1-u)^{n}=e^{n\log(1-u)}$ is concentrated near $u=0$, Laplace's method
(Tome~4, Theorem~11.9; here it is even exact: the integral is $B(\frac32,n+\frac12)$) gives $\int_0^1(1-u)^{n-1/2}\sqrt u\dd u\sim\int_0^\infty e^{-nu}\sqrt u\dd u=\Gamma(\frac32)n^{-3/2}=\frac{\sqrt\pi}2n^{-3/2}$.
So $C_n\sim\frac{4^n}{\sqrt\pi\,n^{3/2}}$.

\tools{Tome~4, Exercise~5.3, Proposition~5.12 ($\Gamma(\frac12)$), Theorem~11.9 (Laplace); Tome~3, Exercise~1.7 (Catalan numbers count Dyck paths).}
\leads (a) is the moment form of Wigner's semicircle law (1955): the eigenvalues of a large random symmetric matrix, scaled by
$\sqrt n$, are distributed with density $\frac1{2\pi}\sqrt{4-x^2}$, and the proof computes $\frac1n\operatorname{tr}(A^{2k})$ by counting closed walks that
turn out to be Dyck paths. (b) is the Marchenko--Pastur law for $AA^{\mathsf T}$. Random matrices return in Part~III.
\end{solution}

\begin{solution}{pb:fracint}
\idea After $x=\frac1t$, the integral runs over $[1,\infty[$ with weight $\frac1{t^2}$, and on each $[k,k+1[$ the fractional part is a
polynomial. The sums that come out are partial sums of the harmonic series and of $\log k$, so $\gamma$ and Stirling's constant appear.

\writeup (a) $\int_0^1\{\frac1x\}\dd x=\int_1^\infty\frac{\{t\}}{t^2}\dd t$, and
\[
\int_k^{k+1}\frac{t-k}{t^2}\dd t=\log\frac{k+1}k-\frac1{k+1}.
\]
Summing for $k\le K$ gives $\log(K+1)-(H_{K+1}-1)\to1-\gamma$ (Tome~4, Corollary~11.6).

(b) Similarly
\[
\int_k^{k+1}\frac{(t-k)^2}{t^2}\dd t=1-2k\log\frac{k+1}k+\frac k{k+1}.
\]
Now $\sum_{k\le K}\frac k{k+1}=K+1-H_{K+1}$ and, by summation by parts,
$\sum_{k\le K}k\log\frac{k+1}k=K\log(K+1)-\log K!$. So the integral over $[1,K+1]$ is
\[
2K+1-H_{K+1}-2K\log(K+1)+2\log K!.
\]
By Stirling (Tome~4, Theorem~5.8), $2\log K!=2K\log K-2K+\log(2\pi K)+o(1)$, while $2K\log(K+1)=2K\log K+2+o(1)$ and $H_{K+1}=\log K+\gamma+o(1)$.
Everything cancels except $1-2-\gamma+\log(2\pi)$: the limit is $\log(2\pi)-\gamma-1=0.2606\ldots$

(c) On $]\frac1{k+1},\frac1k]$, $\lfloor\frac1x\rfloor=k$, so the integral is $\sum_k\frac1k\bigl(\frac1k-\frac1{k+1}\bigr)=\zeta(2)-\sum_k\frac1{k(k+1)}=\zeta(2)-1$.

\tools{Tome~4, Corollary~11.6 (harmonic numbers), Theorem~5.8 (Stirling).}
\leads $\int_0^1\{\frac1x\}x^{s-1}\dd x=\frac1{s-1}-\frac{\zeta(s)}s$ for $\operatorname{Re}s>1$, and the left side is holomorphic for $\operatorname{Re}s>0$; this is how the
fractional part continues $\zeta$ to $\operatorname{Re}s>0$ (Tome~4, Chapter~10). Part (a) is the value at $s=1$, where the pole of $\zeta$ cancels.
\end{solution}

\begin{solution}{pb:fracmean}
\idea $\{\frac xn\}=\frac xn-\lfloor\frac xn\rfloor$, and both sums are classical: the first is $x(\log x+\gamma)$, the second is the
divisor sum $x\log x+(2\gamma-1)x$. The difference keeps only the constants.

\writeup (a) By Tome~5, Lemma~2.5(a), $\sum_{n\le x}\frac xn=x\log x+\gamma x+O(1)$. Counting the pairs $(n,m)$ with $nm\le x$ in two ways,
$\sum_{n\le x}\lfloor\frac xn\rfloor=\sum_{m\le x}\tau(m)=x\log x+(2\gamma-1)x+O(\sqrt x)$ by Dirichlet's theorem (Tome~5, Chapter~2, via the hyperbola
method, Lemma~2.8). Subtract: $\sum_{n\le x}\{\frac xn\}=(1-\gamma)x+O(\sqrt x)$.

(b) $N\bmod n=N-n\lfloor\frac Nn\rfloor$, and $\sum_{n\le N}n\lfloor\frac Nn\rfloor=\sum_{nm\le N}n=\sum_{k\le N}\sigma(k)=\frac{\pi^2}{12}N^2+O(N\log N)$ by Tome~5,
Theorem~2.6(b). So the sum is $N^2-\frac{\pi^2}{12}N^2+O(N\log N)$.

(c) With $t=\frac nx$, $\frac1x\sum_{n\le x}\{\frac xn\}$ is a Riemann sum of $\{\frac1t\}$ on $]0,1]$, whose integral is $1-\gamma$ (Problem~\ref{pb:fracint}). Likewise
$\frac2{N^2}\sum_nn\{\frac Nn\}=\frac2{N^2}\sum(N\bmod n)$ is a Riemann sum of $2t\{\frac1t\}$, and by (b) it tends to $2-\frac{\pi^2}6$; so
$\int_0^12t\{\frac1t\}\dd t=2-\frac{\pi^2}6$, which can also be checked directly as in Problem~\ref{pb:fracint}: $2\int_1^\infty\{u\}u^{-3}\dd u=2-\zeta(2)$.
These Riemann-sum readings are heuristics, since $\{\frac1t\}$ jumps infinitely often near $0$; the error terms come from (a) and (b),
that is, from number theory.

\tools{Tome~5, Lemmas~2.5 and~2.8, Theorem~2.6; Problem~\ref{pb:fracint}.}
\leads The error in (a) is $-\Delta(x)+O(1)$, where $\Delta$ is the error in Dirichlet's divisor problem; so the question ``how regular
are the remainders of $x$ modulo $n$?'' is the divisor problem itself, open since 1849.
\end{solution}

\begin{solution}{pb:quotientlaw}
\idea The event $\lfloor1/U\rfloor=k$ is the interval $]\frac1{k+1},\frac1k]$, and inside it $\{1/U\}=\frac1U-k$. Everything is a sum over $k$ of
lengths of intervals.

\writeup (a) For $u\in[0,1]$, $P(\{1/U\}\le u)=\sum_{k\ge1}P\bigl(k\le\frac1U\le k+u\bigr)=\sum_{k\ge1}\bigl(\frac1k-\frac1{k+u}\bigr)$. The series of derivatives
$\sum\frac1{(k+u)^2}$ converges uniformly on $[0,1]$, so we may differentiate term by term (Tome~4, Chapter~2): the density is $g$. Each term
decreases in $u$; $g(0)=\zeta(2)$ and $g(1)=\zeta(2)-1$.

(b) $P(\lfloor1/U\rfloor\text{ odd})=\sum_{k\text{ odd}}\bigl(\frac1k-\frac1{k+1}\bigr)=1-\frac12+\frac13-\cdots=\log2$. And
$P(\{1/U\}<\frac12)=\sum_k\bigl(\frac1k-\frac1{k+1/2}\bigr)=2\sum_k\bigl(\frac1{2k}-\frac1{2k+1}\bigr)=2(1-\log2)$, because $\sum_k\bigl(\frac1{2k}-\frac1{2k+1}\bigr)=1-(1-\frac12+\frac13-\cdots)$.

(c) $\lfloor x/n\rfloor=k$ if and only if $\frac x{k+1}<n\le\frac xk$, which holds for $x\bigl(\frac1k-\frac1{k+1}\bigr)+O(1)$ integers $n$. Fix $K$. The $n\le x$
with $\lfloor x/n\rfloor\le K$ and odd are $x\sum_{k\le K,\,k\text{ odd}}\bigl(\frac1k-\frac1{k+1}\bigr)+O(K)$ in number, and the other $n$ with $\lfloor x/n\rfloor$ odd
satisfy $n\le\frac x{K+1}$. So the proportion lies, up to $O(\frac Kx)$, between $s_K=\sum_{k\le K\text{ odd}}(\frac1k-\frac1{k+1})$ and $s_K+\frac1{K+1}$. Let
$x\to\infty$, then $K\to\infty$: the proportion tends to $\log2$.

\tools{Tome~3, Chapter~5 (densities); Tome~4, Chapter~2 (term-by-term differentiation); Tome~0, Chapter~5 (alternating harmonic series).}
\leads The law of $\{1/U\}$ is the first step of the Gauss map $x\mapsto\{1/x\}$ of continued fractions; iterating it drives every reasonable
initial law to the Gauss measure $\frac{\mathrm dx}{(1+x)\log2}$ (Problem~\ref{pb:gausskuzmin} of Chapter~6). Note that $\log2$ appears here as the chance that
the first partial quotient of a random number is odd.
\end{solution}

\begin{solution}{pb:primefrac}
\idea The values $\{\frac xp\}$ are $\varphi(\frac px)$ with $\varphi(t)=\{\frac1t\}$. By the prime number theorem the points $\frac px$ are spread over
$]0,1]$ with density $\frac x{\log x}\dd t$ on every fixed interval $[a,b]\subset]0,1]$, because $\log(xt)\sim\log x$ there. So the average of
$\varphi(\frac px)$ should be $\int_0^1\varphi=1-\gamma$. The work is in the approximation.

\writeup Let $\varepsilon>0$ and $K\ge2$. On $[\frac1K,1]$, $\varphi$ is bounded with finitely many discontinuities, hence Riemann integrable: there
are step functions $s\le\varphi\le S$ on $[\frac1K,1]$, constant on the intervals $]t_{j-1},t_j]$ of a subdivision $\frac1K=t_0<\cdots<t_m=1$, with
$\int(S-s)<\varepsilon$. By the prime number theorem (Tome~5, Theorem~5.5), for each fixed $t\in]0,1]$, $\pi(xt)=\frac{xt}{\log(xt)}(1+o(1))=\frac{xt}{\log x}(1+o(1))$.
Hence, for each $j$, the number of primes in $]xt_{j-1},xt_j]$ is $\frac x{\log x}(t_j-t_{j-1})+o(\frac x{\log x})$, and
\[
\sum_{x/K<p\le x}S\Bigl(\frac px\Bigr)=\frac x{\log x}\int_{1/K}^1S+o\Bigl(\frac x{\log x}\Bigr),
\]
and the same for $s$. The sum of $\varphi(\frac px)$ is squeezed between these two. The primes $p\le\frac xK$ contribute between $0$ and
$\pi(\frac xK)\le\frac{2x}{K\log x}$ for large $x$. Finally $\int_0^{1/K}\varphi\le\frac1K$. Dividing by $\pi(x)\sim\frac x{\log x}$,
\[
\limsup_{x\to\infty}\Bigl|\frac1{\pi(x)}\sum_{p\le x}\Bigl\{\frac xp\Bigr\}-\int_0^1\varphi\Bigr|\le\varepsilon+\frac3K,
\]
and $\varepsilon$, $K$ are arbitrary. By Problem~\ref{pb:fracint}(a), $\int_0^1\varphi=1-\gamma$.

\tools{Tome~5, Theorem~5.5 (prime number theorem); Tome~4, Chapter~4 (Riemann integrability, step functions); Problem~\ref{pb:fracint}.}
\leads De la Vallée Poussin proved this in 1898, two years after his proof of the prime number theorem, as an example of a ``mean
value'' over primes. The rate of convergence is tied to the error term in the prime number theorem, and therefore to the zeros of
$\zeta$: the same remark as for the divisor problem, with primes instead of integers.
\end{solution}

\begin{solution}{pb:passone}
\idea $N>n$ means $U_1+\cdots+U_n\le1$, whose probability is the volume of a simplex; and $L\ge n$ means $U_1<\cdots<U_n$, one of $n!$ equally
likely orderings. Both are $\frac1{n!}$, and the partial-sum map shows that this is the same fact.

\writeup (c) first. The linear map
\[
\Phi(u_1,\dots,u_n)=(u_1,\,u_1+u_2,\,\dots,\,u_1+\cdots+u_n)
\]
has a triangular matrix with ones on the diagonal, so it
preserves volume, and it maps $\{u\in[0,1]^n:u_1+\cdots+u_n\le1\}$ bijectively onto $\{0\le s_1\le s_2\le\cdots\le s_n\le1\}$ (the inverse takes
differences, which are non-negative and at most $1$).

(a) Hence $P(N>n)=P(U_1+\cdots+U_n\le1)=\operatorname{vol}\{0\le s_1\le\cdots\le s_n\le1\}$. The $n!$ regions $\{s_{\sigma(1)}\le\cdots\le s_{\sigma(n)}\}$ cover the cube,
overlap in volume $0$ and have the same volume by symmetry, so each has volume $\frac1{n!}$. Then $\mathbb E[N]=\sum_{n\ge0}P(N>n)=\sum\frac1{n!}=e$, and
$\mathbb E[N^2]=\sum_{n\ge0}(2n+1)P(N>n)=2\sum_{n\ge1}\frac1{(n-1)!}+e=3e$, so $\operatorname{Var}N=3e-e^2=0.7657\ldots$

(b) $L\ge n$ if and only if $U_1<\cdots<U_n$, of probability $\frac1{n!}$ (the same symmetry argument). So $P(L+1>n)=P(L\ge n)=\frac1{n!}=P(N>n)$ for
every $n$.

\tools{Tome~3, Chapter~5 (uniform variables, $\mathbb E[X]=\sum P(X>n)$); Tome~0, Chapter~8.}
\leads Bijections like $\Phi$, which turn simplices into ordered tuples, are behind Stanley's proof in the next problem and behind the
volume computations of polytopes in combinatorics (Ehrhart theory).
\end{solution}

\begin{solution}{pb:irwinhall}
\idea Reading $S_i$ modulo $1$ turns the walk $S_0,S_1,\dots$ into a sequence of independent uniform points on a circle, and $\lfloor S_n\rfloor$
counts the times the walk crosses an integer, that is, the times the point on the circle wraps around: the descents.

\writeup (a) $P(S_n\le x)$ is the volume of $Q(x)=\{u\in[0,1]^n:\sum u_i\le x\}$. For $J\subseteq[n]$, the set $\{u\in[0,\infty[^n:\sum u_i\le x,\ u_i\ge1\text{ for }i\in J\}$ is a
translate of $\{v\ge0:\sum v_i\le x-|J|\}$, of volume $\frac{(x-|J|)_+^n}{n!}$ by Problem~\ref{pb:passone}. By inclusion--exclusion on the set of coordinates
exceeding $1$,
\[
\operatorname{vol}Q(x)=\sum_{J\subseteq[n]}(-1)^{|J|}\frac{(x-|J|)_+^n}{n!}=\frac1{n!}\sum_j(-1)^j\binom nj(x-j)_+^n.
\]

(b) $V_i=\{V_{i-1}+U_i\}$. Given $V_1,\dots,V_{i-1}$, the variable $U_i$ is still uniform and independent of them, and $u\mapsto\{v+u\}$ maps the
uniform law on $[0,1[$ to itself (a rotation of the circle $\R/\Z$). So $V_i$ is uniform and independent of $(V_1,\dots,V_{i-1})$: the $V_i$
are independent uniform. As $0\le U_i<1$, $\lfloor S_i\rfloor-\lfloor S_{i-1}\rfloor\in\{0,1\}$, and it equals $1$ exactly when $V_i<V_{i-1}$. Since $V_1=U_1\ge V_0=0$,
$\lfloor S_n\rfloor$ is the number of $i\in\{1,\dots,n-1\}$ with $V_{i+1}<V_i$.

(c) The $V_i$ are almost surely distinct, and their relative order is a uniform random permutation $\sigma$ (Tome~3, Chapter~5). By (b),
$\lfloor S_n\rfloor$ is the number of descents of $\sigma$, so $P(\lfloor S_n\rfloor=k)=\frac{A(n,k)}{n!}$. With $F(x)=n!\,P(S_n\le x)$ from (a),
\[
A(n,k)=F(k+1)-F(k)=\sum_j(-1)^j\binom nj(k+1-j)^n_+-\sum_j(-1)^j\binom nj(k-j)_+^n.
\]
Shift $j$ to $j-1$ in the second sum and use $\binom nj+\binom n{j-1}=\binom{n+1}j$: $A(n,k)=\sum_{j=0}^k(-1)^j\binom{n+1}j(k+1-j)^n$. For $n=4$ this gives
$1,11,11,1$ (Figure~\ref{fig:irwinhall}).

\tools{Problem~\ref{pb:passone}; Tome~0, Chapter~8 (inclusion--exclusion); Tome~3, Chapter~5.}
\leads Richard Stanley found (b) in 1977. Eulerian numbers already appeared in Tome~3, Chapter~8, counting rising sequences in the
analysis of riffle shuffles, where the Bayer--Diaconis formula has exactly this flavour. Since $S_n$ is approximately normal, so is the
number of descents of a random permutation: mean $\frac{n-1}2$, variance $\frac{n+1}{12}$.
\end{solution}

\begin{solution}{pb:sincpower}
\idea $\frac{\sin t}t=\mathbb E[e^{itV}]$ for $V$ uniform on $[-1,1]$. So $(\frac{\sin t}t)^n$ is the characteristic function of a sum of $n$ such
variables, and the inversion formula at $0$ turns the integral into a value of the density, which Problem~\ref{pb:irwinhall} computes.

\writeup (a) $\mathbb E[e^{itV}]=\frac12\int_{-1}^1e^{itv}\dd v=\frac{\sin t}t$. Let $Y=V_1+\cdots+V_n$; $\varphi_Y(t)=(\frac{\sin t}t)^n$ is integrable for $n\ge2$, so by
Tome~3, Theorem~6.8(b), $Y$ has a continuous density $f$ with $f(0)=\frac1{2\pi}\int(\frac{\sin t}t)^n\dd t$. Now $V_i=2U_i-1$ with $U_i$ uniform on $[0,1]$,
so $Y=2S_n-n$ and $f(0)=\frac12f_{S_n}(\frac n2)$, where $f_{S_n}(x)=\frac1{(n-1)!}\sum_j(-1)^j\binom nj(x-j)_+^{n-1}$ is the derivative of Problem~\ref{pb:irwinhall}(a). So
\[
\int\Bigl(\frac{\sin t}t\Bigr)^n\dd t=\pi f_{S_n}\Bigl(\frac n2\Bigr)=\frac\pi{(n-1)!}\sum_{j<n/2}(-1)^j\binom nj\Bigl(\frac n2-j\Bigr)^{n-1},
\]
which is the formula. For $n=2,3,4$: $\frac\pi{2}\cdot2=\pi$, $\frac\pi{8}(9-3)=\frac{3\pi}4$, $\frac\pi{48}(64-32)=\frac{2\pi}3$.

(b) On $[-\pi,\pi]$, $\frac{\sin t}t>0$ except at $\pm\pi$, and $\phi(t)=\log\frac{\sin t}t$ is smooth on $]-\pi,\pi[$ with a unique maximum $\phi(0)=0$ and
$\phi(t)=-\frac{t^2}6+O(t^4)$. By Laplace's method (Tome~4, Theorem~11.9), $\int_{-\pi}^\pi e^{n\phi}\sim\sqrt{\frac{2\pi}{n/3}}=\sqrt{\frac{6\pi}n}$. Outside,
$\bigl|\int_{|t|>\pi}(\frac{\sin t}t)^n\bigr|\le2\int_\pi^\infty t^{-n}\dd t=\frac{2\pi^{1-n}}{n-1}$, which is negligible.

\tools{Tome~3, Theorem~6.8 (inversion); Tome~4, Theorem~11.9 (Laplace's method); Problem~\ref{pb:irwinhall}.}
\leads (b) says $f(0)\sim\frac1{\sqrt{2\pi n/3}}$, the value at $0$ of the normal density of variance $\frac n3$: a local central limit theorem.
Integrals of this kind measure sections of the cube; with them Hensley (1979) and Ball (1986) proved that every central hyperplane
section of the unit cube of $\R^n$ has volume between $1$ and $\sqrt2$.
\end{solution}

\begin{solution}{pb:borwein}
\idea The product of the sinc factors with $k\ge1$ is the characteristic function of $Y$. The remaining factor $\operatorname{sinc}(a_0x)$ is the
Fourier transform of the indicator of $[-a_0,a_0]$, so the integral asks how often $|Y|<a_0$. As long as $\sum_{k\ge1}a_k\le a_0$, the answer
is always.

\writeup (a) Let $\operatorname{Si}(z)=\int_0^z\frac{\sin u}u\dd u$; it is bounded on $\R$ and tends to $\pm\frac\pi2$ as $z\to\pm\infty$ (Tome~4, Theorem~4.11). For real $y$
and $R>0$,
\[
\begin{aligned}\int_{-R}^R\operatorname{sinc}(a_0x)\cos(xy)\dd x&=\frac1{a_0}\int_0^R\frac{\sin((a_0+y)x)+\sin((a_0-y)x)}x\dd x\\&=\frac{\operatorname{Si}((a_0+y)R)+\operatorname{Si}((a_0-y)R)}{a_0},\end{aligned}
\]
which is bounded uniformly in $R$ and $y$, and tends to $\frac\pi{a_0}$ if $|y|<a_0$ and to $0$ if $|y|>a_0$. Since $\mathbb E[e^{ixY}]=\prod_{k\ge1}\operatorname{sinc}(a_kx)$
is real and even, Fubini on the bounded set $[-R,R]$ gives
\[
\int_{-R}^R\prod_{k=0}^m\operatorname{sinc}(a_kx)\dd x=\mathbb E\Bigl[\int_{-R}^R\operatorname{sinc}(a_0x)\cos(xY)\dd x\Bigr].
\]
Let $R\to\infty$. On the left, the integrand is $O(x^{-2})$ (as $m\ge1$), so the integral converges to $2\int_0^\infty$. On the right,
bounded convergence (Tome~3, Chapter~5) gives $\frac\pi{a_0}P(|Y|<a_0)$, as $Y$ has a density and $P(|Y|=a_0)=0$.

(b) With $a_k=\frac1{2k+1}$, $|Y|\le\sum_{k=1}^ma_k$, and $\frac13+\frac15+\cdots+\frac1{13}=0.955\ldots<1$, so $P(|Y|<1)=1$ for $m\le6$ (and $m=0$ is
Dirichlet's integral itself). For $m=7$, $\sum_{k=1}^7a_k=1+\delta$ with $\delta=\frac{982}{45045}>0$, and $Y$ exceeds $1$ with positive probability: all
$V_k$ close to $1$.

(c) By symmetry $P(|Y|\ge1)=2P(Y>1)$. Put $w_k=1-V_k\in[0,2]$: $Y>1$ means $\sum_ka_kw_k<\delta$. As $\delta<2a_k$ for every $k$, this set of
$w$ is the full simplex $\{w\ge0,\sum a_kw_k<\delta\}$, of volume $\frac{\delta^7}{7!\,a_1\cdots a_7}$, and the density of $(w_k)$ is $2^{-7}$. Hence
\[
\frac\pi2-\int_0^\infty\prod_{k=0}^7\operatorname{sinc}\frac x{2k+1}\dd x=\frac\pi2\cdot\frac{2\delta^7}{2^7\,7!\,a_1\cdots a_7}=\frac{\pi\,\delta^7\cdot3\cdot5\cdots15}{2^7\cdot7!}=2.310\ldots\cdot10^{-11}.
\]

\tools{Tome~4, Theorem~4.11 (Dirichlet's integral); Tome~3, Chapter~5 (bounded convergence) and Chapter~6 (characteristic functions).}
\leads D.~Borwein and J.~M.~Borwein (2001) found the pattern and its failure; the argument of (a) explains every integral of this
family at once. Baillie, Borwein and Borwein (2008) showed that in many such cases the sum $\sum_{n\ge1}\prod_k\operatorname{sinc}(a_kn)$ equals the
integral minus $\frac12$, a case of Poisson summation (Tome~4, Theorem~8.9): the Fourier transform of the product vanishes outside
$[-\sum a_k,\sum a_k]$.
\end{solution}

\begin{solution}{pb:stick}
\idea On a circle, $n$ independent uniform points are all alike, so the $n$ arcs are exchangeable; and an arc of length at least $\frac12$
is rare and unique. For (b), a random point falls into a piece with probability equal to its length: long pieces are favoured.

\writeup (a) Take $n$ independent uniform points $P_1,\dots,P_n$ on a circle of length $1$, and cut the circle at $P_1$: the other points become
$n-1$ independent uniform points on a stick, and the arcs become the pieces. Let $A_j$ be the arc from $P_j$ clockwise to the next point.
Two arcs cannot both have length $\ge\frac12$ except on a null event, and $A_j\ge\frac12$ means that the $n-1$ other points avoid the half circle
following $P_j$, of probability $2^{-(n-1)}$. So $P(\max_jA_j\ge\frac12)=n2^{-(n-1)}$. (Pieces $\ell_1,\dots,\ell_n$ bound an $n$-gon if and only if the
largest is less than the sum of the others, that is, less than $\frac12$.)

(b) Given the pieces, $T$ falls in the $i$-th one with probability $L_i$, so the expected length is $\mathbb E[\sum_iL_i^2]=n\mathbb E[L_1^2]$. Now
$P(L_1>t)=(1-t)^{n-1}$ (all cuts after $t$), so $\mathbb E[L_1^2]=\int_0^12t(1-t)^{n-1}\dd t=\frac2{n(n+1)}$. Hence $\frac2{n+1}$.

\tools{Tome~3, Chapter~5 (uniform variables, $\mathbb E[X^2]=\int2tP(X>t)\dd t$).}
\leads (b) is the inspection paradox, the reason why the bus you wait for seems to come from a long gap (Tome~3, Chapter~9). Its cousin for
the Poisson process says that the gap containing a fixed time has twice the mean length.
\end{solution}

\begin{solution}{pb:spacings}
\idea Uniform spacings are normalised exponentials: this Rényi--Lévy representation turns maxima of spacings into maxima of independent
variables, which are easy.

\writeup (a) The joint density of $(E_1,\dots,E_n)$ is $e^{-(e_1+\cdots+e_n)}$ on $]0,\infty[^n$. Change variables to $s=\sum e_i$ and $\ell_i=\frac{e_i}s$ for $i<n$
(and $\ell_n=1-\sum_{i<n}\ell_i$); the inverse $e_i=s\ell_i$ ($i<n$), $e_n=s(1-\sum\ell_i)$ has Jacobian $s^{n-1}$. So $(S,\ell_1,\dots,\ell_{n-1})$ has density
$s^{n-1}e^{-s}$ times the indicator of the simplex $\Delta=\{\ell_i>0,\sum_{i<n}\ell_i<1\}$: $S$ is independent of $(\ell_i)$, which is uniform on $\Delta$ with
density $(n-1)!$. On the other side, the order statistics of $n-1$ uniforms have density $(n-1)!$ on $\{0<t_1<\cdots<t_{n-1}<1\}$ (Tome~3, Chapter~9), and
the differences $L_i=t_i-t_{i-1}$ ($i<n$) are obtained by a linear map of determinant $1$ onto $\Delta$: the same law.

(b) $P(\max E_i\le t)=(1-e^{-t})^n$, so with $v=1-e^{-t}$,
\[
\mathbb E[\max E_i]=\int_0^\infty\bigl(1-(1-e^{-t})^n\bigr)\dd t=\int_0^1\frac{1-v^n}{1-v}\dd v=H_n.
\]
By (a), $\max E_i$ has the law of $S\cdot M$ with $M=\max L_i$ independent of $S$, so $H_n=\mathbb E[S]\,\mathbb E[M]=n\,\mathbb E[M]$. Similarly $\min E_i$ is
exponential of parameter $n$ ($P(\min>t)=e^{-nt}$), of mean $\frac1n=n\,\mathbb E[\min L_i]$.

\tools{Tome~3, Chapter~5 (change of variables, Exercise~5.8 for $n=2$) and Chapter~9 (order statistics of uniforms).}
\leads So the longest piece is about $\frac{\log n}n$, and Lévy (1939) showed that $n\max L_i-\log n$ converges in law to the Gumbel distribution.
The question ``how large is the largest gap?'' returns for primes: Cramér's model predicts gaps of size $(\log p)^2$, which Chapter~10 of
Tome~5 discusses.
\end{solution}

\begin{solution}{pb:sophomore}
\idea $x^{-x}=e^{-x\log x}$, and each term of the exponential series integrates to a Gamma integral; the $k!$ cancels.

\writeup (a) $x^{-x}=\sum_{k\ge0}\frac{(-x\log x)^k}{k!}$ with non-negative terms on $]0,1]$, so we may integrate term by term. With $x=e^{-u}$,
\[
\int_0^1x^k(-\log x)^k\dd x=\int_0^\infty u^ke^{-(k+1)u}\dd u=\frac{k!}{(k+1)^{k+1}},
\]
so $\int_0^1x^{-x}\dd x=\sum_{k\ge0}\frac1{(k+1)^{k+1}}$. For $x^x$ the terms have alternating signs, but their absolute values are those of $x^{-x}$,
whose integral converges, so dominated convergence applies and gives $\sum_{k\ge0}\frac{(-1)^k}{(k+1)^{k+1}}$.

(b) The tail after $n=N$ is at most $\sum_{n>N}(N+1)^{-n}\le\frac{(N+1)^{-N-1}}{1-1/(N+1)}$. For $N=9$ this is $1.1\cdot10^{-10}$, too large to guarantee ten
decimals; for $N=10$ it is $3.9\cdot10^{-12}$. Ten terms: $1.2912859970\ldots$

\tools{Tome~4, Chapter~5 (Gamma integral) and Chapter~4 (term-by-term integration).}
\leads Johann Bernoulli found these identities in 1697. The name ``sophomore's dream'' was coined by J.~Borwein, D.~Bailey and
R.~Girgensohn (2004), for an identity that looks like a student's mistake and is true.
\end{solution}

\section*{Map of the chapter}

\begin{center}
\small
\adjustbox{max width=\linewidth}{\begin{tabular}{lll}
Problem & Topic & Tomes used\\\hline
\ref{pb:logtwo} & $\log2$ is irrational & 0, 1, 4\\
\ref{pb:hermite} & $e^r$ and $\log r$ are irrational & 0, 4\\
\ref{pb:doubleint} & double integrals for $\zeta(2)$, $\zeta(3)$ & 4\\
\ref{pb:beukerstwo} & Beukers for $\zeta(2)$ & 0, 1, 4, 5\\
\ref{pb:beukersthree} & Beukers for $\zeta(3)$ & 1, 4, 5\\
\ref{pb:sondow} & $\gamma$ and $\log\frac4\pi$ as integrals & 4\\
\ref{pb:dalzell} & $\pi<\frac{22}7$ & 2, 4\\
\ref{pb:centralsum} & $\sum1/\binom{2n}n$ & 0, 4\\
\ref{pb:arcsin} & $\arcsin^2$ and $\frac{\pi^2}{18}$ & 0, 4\\
\ref{pb:lehmer} & series with $\binom{2n}n4^{-n}$ & 0, 4\\
\ref{pb:catalanmoments} & Catalan numbers as moments & 3, 4\\
\ref{pb:fracint} & integrals of $\{1/x\}$ & 4\\
\ref{pb:fracmean} & remainders on average & 5\\
\ref{pb:quotientlaw} & the law of $\{1/U\}$ & 0, 3, 4\\
\ref{pb:primefrac} & remainders modulo primes & 4, 5\\
\ref{pb:passone} & waiting to pass $1$ & 0, 3\\
\ref{pb:irwinhall} & sums of uniforms, Eulerian numbers & 0, 3\\
\ref{pb:sincpower} & powers of $\frac{\sin t}t$ & 3, 4\\
\ref{pb:borwein} & the Borwein integrals & 3, 4\\
\ref{pb:stick} & breaking a stick & 3\\
\ref{pb:spacings} & longest and shortest pieces & 3\\
\ref{pb:sophomore} & the sophomore's dream & 4\\
\end{tabular}}
\end{center}

\section*{Sources and further reading}

\begin{itemize}
\item F.~Beukers, A note on the irrationality of $\zeta(2)$ and $\zeta(3)$, \emph{Bulletin of the London Mathematical Society} 11 (1979),
268--272; F.~Beukers, \emph{A passion for numbers}, farewell lecture, Utrecht University, 14 November 2019 (the epigraph, and the account
of the Marseille lecture quoted in the opening).
\item A.~van der Poorten, A proof that Euler missed\dots{} Apéry's proof of the irrationality of $\zeta(3)$, \emph{Mathematical Intelligencer}
1 (1979), 195--203 (the quotation in the opening).
\item R.~Apéry, Irrationalité de $\zeta(2)$ et $\zeta(3)$, \emph{Astérisque} 61 (1979), 11--13.
\item T.~Rivoal, La fonction zêta de Riemann prend une infinité de valeurs irrationnelles aux entiers impairs, \emph{Comptes Rendus de
l'Académie des Sciences, Série I} 331 (2000), 267--270; W.~Zudilin, One of the numbers $\zeta(5),\zeta(7),\zeta(9),\zeta(11)$ is irrational,
\emph{Russian Mathematical Surveys} 56 (2001), 774--776.
\item R.~Marcovecchio, The Rhin--Viola method for $\log2$, \emph{Acta Arithmetica} 139 (2009), 147--184; G.~Rhin and C.~Viola, The group
structure for $\zeta(3)$, \emph{Acta Arithmetica} 97 (2001), 269--293.
\item J.~Sondow, Double integrals for Euler's constant and $\ln\frac4\pi$ and an analog of Hadjicostas's formula, \emph{American
Mathematical Monthly} 112 (2005), 61--65; J.~Sondow, Criteria for irrationality of Euler's constant, \emph{Proceedings of the AMS} 131 (2003),
3335--3344.
\item D.~P.~Dalzell, On $22/7$, \emph{Journal of the London Mathematical Society} 19 (1944), 133--134; S.~K.~Lucas, Approximations to $\pi$
derived from integrals with nonnegative integrands, \emph{American Mathematical Monthly} 116 (2009), 166--172; D.~Zeilberger and
W.~Zudilin, The irrationality measure of $\pi$ is at most $7.103205334137\ldots$, \emph{Moscow Journal of Combinatorics and Number
Theory} 9 (2020), 407--419.
\item D.~H.~Lehmer, Interesting series involving the central binomial coefficient, \emph{American Mathematical Monthly} 92 (1985),
449--457; L.~Comtet, \emph{Advanced Combinatorics}, Reidel, 1974; J.~M.~Borwein and D.~M.~Bradley, Empirically determined Apéry-like
formulae for $\zeta(4n+3)$, \emph{Experimental Mathematics} 6 (1997), 181--194.
\item E.~P.~Wigner, Characteristic vectors of bordered matrices with infinite dimensions, \emph{Annals of Mathematics} 62 (1955), 548--564.
\item C.~de la Vallée Poussin, Sur les valeurs moyennes de certaines fonctions arithmétiques, \emph{Annales de la Société Scientifique de
Bruxelles} 22 (1898), 84--90.
\item R.~P.~Stanley, Eulerian partitions of a unit hypercube, in \emph{Higher Combinatorics} (M.~Aigner, ed.), Reidel, 1977, p.~49.
\item D.~Borwein and J.~M.~Borwein, Some remarkable properties of sinc and related integrals, \emph{The Ramanujan Journal} 5 (2001), 73--89;
R.~Baillie, D.~Borwein and J.~M.~Borwein, Surprising sinc sums and integrals, \emph{American Mathematical Monthly} 115 (2008), 888--901.
\item K.~Ball, Cube slicing in $\R^n$, \emph{Proceedings of the AMS} 97 (1986), 465--473.
\item P.~Lévy, Sur la division d'un segment par des points choisis au hasard, \emph{Comptes Rendus de l'Académie des Sciences} 208 (1939),
147--149.
\item J.~Borwein, D.~Bailey and R.~Girgensohn, \emph{Experimentation in Mathematics}, A K Peters, 2004 (the sophomore's dream, Apéry-like
formulas, the Borwein integrals).
\end{itemize}
